\clearpage
\VerdierPhysicalPageAnchor{0244}%
\begingroup\small\setlength{\parskip}{0pt}\setlength{\emergencystretch}{1.5em}
\setlength{\abovedisplayskip}{4pt plus 2pt minus 2pt}%
\setlength{\belowdisplayskip}{4pt plus 2pt minus 2pt}%

\VerdierSemanticAnchor{ch4d-4-5-4-proof-e-exact-diagram}
これにより、各行が完全である可換図式を得る：
\[
\begin{tikzcd}[column sep=2.7em,row sep=2.8em]
F_1^{p+r-1}(Y) \arrow[r] \arrow[d] &
F_1^p(L^\bullet) \arrow[r] \arrow[d,"{m(L^\bullet[p])}"'] &
F_1^p(X) \arrow[r] \arrow[d,"{m(X[p])}"'] &
F_1^{p+r}(Y) \arrow[d] \\
F_2^{p+r-1}(Y) \arrow[r] &
F_2^p(L^\bullet) \arrow[r] &
F_2^p(X) \arrow[r] &
F_2^{p+r}(Y).
\end{tikzcd}
\]
整数 $r$ は、$i=1,2$ および $q\geq p+r-1$ に対して対象 $F_i^q(Y)$ が
零となるように選んだ。したがって、すべての成分が $M$ の元である
$\operatorname{comp}^{b}(\mathcal A)$ の任意の対象 $L^\bullet$ に対して、
$m(L^\bullet)$ が同型であることを示せばよい。

\VerdierSemanticAnchor{ch4d-4-5-4-proof-e-spectral-sequence}
一方、射 $m$ は、複体 $L^\bullet$ に関する関手 $F_1$ と $F_2$ の第一
スペクトル系列の間の射を誘導する。初期項上で、この射
\[
m_1^{p,q}(L^\bullet)\colon F_1^q(L^p)\longrightarrow F_2^q(L^p),
\]
は仮定により同型である。これらのスペクトル系列は定常なので (4.5.3)、
射 $m$ は収束先上の同型を誘導し（第 II 章、4.4.5）、したがって (4.4.2)
により射 $m(L^\bullet)$ は同型である。

\VerdierSemanticAnchor{ch4d-4-5-4-proof-end}
これで主張 e) の証明が完了する。主張 f) は反対圏へ移ることにより主張 e)
から従う。定理が示された。

\VerdierSemanticAnchor{ch4d-4-6}
\medskip\noindent\textbf{4.6. Cartan--Eilenberg のスペクトル対象。}

\VerdierSemanticAnchor{ch4d-4-6-1}
\noindent\textbf{定義 4.6.1.}\quad
$\mathcal A$ の複体 $I^\bullet$ が\emph{Cartan--Eilenberg 入射複体}である
とは、各整数 $i$ に対して対象 $\operatorname{Im}(d_{I^\bullet}^{i})$ と
$H^i(I^\bullet)$ が入射対象であることをいう。

\VerdierSemanticAnchor{ch4d-4-6-2}
\noindent\textbf{4.6.2.}\quad
Cartan--Eilenberg 入射複体 $I^\bullet$ は入射型であることに注意する
(1.1.5)。すなわち、その成分は $\mathcal A$ の入射対象であり、各整数
$i$ に対して対象 $\operatorname{Ker}(d_{I^\bullet}^{i})$ は
$\mathcal A$ の入射対象である。これらの対象を結ぶ三項完全列を用いれば、
このことは直ちに確かめられる。

\VerdierSemanticAnchor{ch4d-4-6-3-a}
\medskip\noindent\textbf{命題 4.6.3.}\quad
\textit{$I^\bullet$ を Cartan--Eilenberg 入射複体とする。}

\noindent\textit{a) 任意の複体 $Y$ に対し、標準射}
\[
\operatorname{Hom}_{K(\mathcal A)}(Y,I^\bullet)
\longrightarrow
\prod_{p\in\mathbb Z}
\operatorname{Hom}_{\mathcal A}\bigl(H^p(Y),H^p(I^\bullet)\bigr)
\]
\textit{は同型である。}
\endgroup
